\pagestyle{plain} % No headers, just page numbers
\pagenumbering{arabic} % Arabic numerals
\setcounter{page}{1}

%%%%
% - what is DA
% - why is it important
% - multiparty
%     much more realistic, but more challenging

    
% - spk interactions/roles are important
%     hidden context

% - LLMs still struggle with it 
% - situational context
    % CF, Structure
    
%%%%

\chapter{\uppercase {Introduction}}
\textit{The success of communication depends on the addressee's ability to infer the communicator's intentions. - Sperber \& Wilson, 1986}. Human communication relies not merely on the explicit semantic content of utterances but critically depends on understanding the underlying intentions that motivate those utterances ~\cite{Austin1962, Searle1969, Grice1975}. Intent understanding serves as the foundation for successful communication, enabling conversational partners to coordinate actions, resolve ambiguities, and build shared understanding~\cite{Clark1996}. Research in pragmatics has long established that the same linguistic content can serve vastly different communicative functions depending on the speaker's goal~\cite{Levinson1983}. similarly, psychological studies demonstrate that misunderstanding intent leads to communication breakdowns, interpersonal conflicts, and task failures~\cite{Thomas1983, Keysar1994}. Importantly, intentions are not always benevolent and speakers may pursue deceptive or manipulative goals, making it equally crucial to recognize malicious intent~\cite{Ekman2001, DePaulo2003, Vrij2008}. Together, these insights underscore a critical reality: effective communication hinges on accurate intent recognition.

As conversational AI systems and Large Language Models (LLMs) become increasingly integrated into daily life, their ability to understand human intent becomes even more crucial for utility and safety. As conversational AI systems increasingly mediate critical interactions in customer service~\cite{chatbots_AI}, personal assistance~\cite{VoiceAssistants}, healthcare~\cite{ConversationalHealthcare}, and education~\cite{ChatbotsEducation}, failures in intent understanding can have serious real-world consequences. In business contexts, intent misclassification in automated customer service systems directly impacts customer satisfaction, retention, and revenue~\cite{Chung2018}. Moreover, the stakes extend beyond individual interactions to societal impact. Research in human-computer interaction demonstrate that repeated intent misclassification erodes user trust in AI systems, creating barriers to adoption even when the technology could provide genuine benefits~\cite{Khadpe2020, Bansal2019}. In collaborative work environments, AI systems that misunderstand human intent disrupt workflow efficiency and team coordination~\cite{Seeber2020}. Altogether, these problems motivate a focused research effort on intent understanding that is robust across domains, interaction modalities, and deployment contexts. The overarching goal of the proposed work is to address this challenge by developing and evaluating a suite of computational models for robust intent understanding, aiming to make AI agents more reliable, adaptable, and safe partners in communication.

Dialog acts (DAs) are the minimal units of communicative intent, characterizing utterances by their conversational function rather than grammatical form and enabling dialogue systems to track state, infer user goals, and respond appropriately~\cite{Searle1969, Bunt2000, Stolcke2000}. Robust DA classification is therefore fundamental to building conversational AI systems that can maintain dialogue state, anticipate user goals, and produce contextually appropriate responses. However, achieving accurate DA understanding in real-time spoken interactions remains challenging because online systems must operate on noisy ASR transcripts and cannot use \textit{oracle} text, while crucial disambiguating cues reside in the audio modality. To address these challenges, we propose a Hierarchical Fusion framework for online multimodal DA classification~\cite{MiahDA} that integrates ASR-generated text with raw audio representations via both early and late fusion stages. Whisper-derived features provide robust acoustic and prosodic features, which are hierarchically combined with textual encodings to preserve fine-grained local cues while enforcing global conversational coherence. This end-to-end design yields state-of-the-art F1 performance on two DA datasets, MRDA and EMOTYDA, demonstrating that principled multimodal fusion substantially improves online DA understanding under realistic ASR noise.

As human-AI interactions grow in complexity, particularly in high-stakes industrial domains like healthcare or finance, the cost of intent misclassification rises sharply. When an AI agent misunderstands a user's intent due to ASR errors, background noise, or unexpected user behaviors, it can lead to irrelevant responses, user frustration, and a complete breakdown of the conversation~\cite{higashinaka-etal-2016-dialogue}. Detecting these Dialogue Breakdowns (DB) in real-time is therefore critical for system reliability. An early DB detection allows the system to take corrective action, such as asking for clarification or escalating to a human operator, before the task fails entirely. In this work, we address this problem by developing a multimodal contextual dialogue breakdown (MultConDB) model~\cite{MiahDB}. The proposed model moves beyond prior text-only methods by incorporating both acoustic signals and textual features, recognizing that many breakdown causes have multimodal signatures. Our model significantly outperforms existing methods in a large-scale industrial setting, proving that multimodal context is essential for maintaining conversational integrity.

While benign intent understanding focuses on accurately interpreting legitimate communicative goals, a complete picture of intent recognition must also address malicious intentions and deception. Deception detection represents another critical facet of intent understanding, with applications spanning criminal investigations, security screening, fraud prevention, and AI safety~\cite{depaulo-54, vrij2008detecting}. Automated deception detection, however, is an extremely difficult task. Humans themselves perform only slightly better than chance, with an average accuracy of around 54\%~\cite{depaulo-54}. Given the human limitations, there is significant interest in whether modern computational models, particularly LLMs, can learn to identify the subtle linguistic and multimodal cues of deception more reliably. To investigate this, we conducted a rigorous and large-scale evaluation of Large Language Models and Large Multimodal Models capabilities for deception detection~\cite{MiahDC}. Using three distinct datasets (RLTD, MU3D, and OpSpam) that cover legal testimonies, interpersonal deception, and deceptive online reviews, we benchmarked various models and prompting strategies. Our findings reveal that while fine-tuned LLMs can achieve state-of-the-art performance on purely textual deception, current LMMs struggle significantly to leverage multimodal cues, highlighting a critical gap in their `understanding' of non-verbal deceptive behaviors.

The exploration of malicious intent leads directly to one of the most pressing, modern challenges in AI safety: LLM jailbreaking. A jailbreak is a form of adversarial attack where a user crafts a prompt to bypass a model's safety filters and elicit harmful, disallowed, or toxic content~\cite{zou2023JB1, wei2023JB2}. The defense against these attacks is contigent upon the ability of the LLMs' malicious intent detection. Current defenses against such attacks are often inadequate. Many rely on external content classifiers, which add latency and can be bypassed by novel attacks, or require multiple LLM calls, which is computationally expensive~\cite{mazeika2024JB3, inan2023JB4}. Our future work will focus on improving the LLM's intrinsic ability to understand and defend against these malicious intents. We propose a novel defense inspired by human strategies for detecting and resisting deception. By training LLMs to think before refusing in human like manner and incorporating real-world deception detection strategies, we aim to develop more robust defenses that generalize across diverse attack vectors.

The research presented in this work collectively addresses the multifaceted challenge of understanding human intent in conversational AI systems. Dialog act classification establishes the foundation by enabling accurate recognition of communicative goals in multimodal interactions. Dialogue breakdown detection builds upon this by identifying when intent understanding fails, allowing systems to recover gracefully. Deception detection extends the framework to malicious intentions, recognizing when speakers deliberately mislead. Finally, jailbreak defense connects these insights by treating adversarial attacks as sophisticated intent manipulation requiring human-inspired detection strategies. This integrated perspective has implications for building conversational AI systems that are not only more capable at understanding what users want to accomplish, but also more reliable in detecting when communication breaks down, more discerning in recognizing deceptive behavior, and more resilient against adversarial manipulation.

